\section{Linear programming and duality}\label{app:C.12}

Every rate of the book is linear in the dark-matter velocity distribution: double the number of particles moving at one velocity, and their contribution to every bin doubles. A limit that should hold for every velocity distribution, the problem of \cref{sec:21.2}, is then the largest value of a linear quantity over all nonnegative distributions that keep the predicted counts below the observed bounds. That is a linear program with infinitely many variables. \Cref{sec:21.2} uses three of its properties: any feasible dual solution certifies an upper bound, the optimum is attained by a halo made of a few streams, and a grid solution errs by an amount that a Lipschitz constant bounds.

\subsection{Finite linear programs}

A linear program in standard form maximizes $\mathbf{c}^{\mathsf{T}}\mathbf{x}$ over vectors $\mathbf{x} \in \mathbb{R}^M$ with $\mathbf{x} \ge 0$ and $\mathbf{A}\mathbf{x} \le \mathbf{b}$, where $\mathbf{A}$ is an $n\times M$ matrix and inequalities between vectors hold componentwise. Its \emph{dual} minimizes $\mathbf{b}^{\mathsf{T}}\boldsymbol{\lambda}$ over $\boldsymbol{\lambda} \in \mathbb{R}^n$ with $\boldsymbol{\lambda} \ge 0$ and $\mathbf{A}^{\mathsf{T}}\boldsymbol{\lambda} \ge \mathbf{c}$. For any feasible pair,
\begin{equation}
  \mathbf{c}^{\mathsf{T}}\mathbf{x} \le (\mathbf{A}^{\mathsf{T}}\boldsymbol{\lambda})^{\mathsf{T}}\mathbf{x} = \boldsymbol{\lambda}^{\mathsf{T}}(\mathbf{A}\mathbf{x}) \le \boldsymbol{\lambda}^{\mathsf{T}}\mathbf{b} .
  \label{eq:C.12-weak}
\end{equation}
The first inequality uses $\mathbf{x} \ge 0$ and $\mathbf{A}^{\mathsf{T}}\boldsymbol{\lambda} \ge \mathbf{c}$, the second $\boldsymbol{\lambda} \ge 0$ and $\mathbf{A}\mathbf{x} \le \mathbf{b}$. This is \emph{weak duality}: every dual-feasible $\boldsymbol{\lambda}$ bounds the primal maximum from above, whether or not anyone has found the maximum. \emph{Strong duality}, the equality of the two optima, holds for finite programs whenever both are feasible.

The optimum of a finite program is attained at a point with at most $n$ nonzero components, one per constraint. Suppose an optimal $\mathbf{x}$ has more than $n$ positive components. The corresponding columns of $\mathbf{A}$, more than $n$ vectors in $\mathbb{R}^n$, are linearly dependent, so there is a nonzero $\mathbf{d}$, supported on those components, with $\mathbf{A}\mathbf{d} = 0$. Then $\mathbf{x} \pm t\mathbf{d}$ satisfies every constraint for small $t > 0$, and since $\mathbf{x}$ is optimal, $\mathbf{c}^{\mathsf{T}}\mathbf{d} = 0$. Increasing $t$ until the first component of $\mathbf{x} + t\mathbf{d}$ or $\mathbf{x} - t\mathbf{d}$ reaches zero yields another optimal point with one positive component fewer. Repeating the step leaves at most $n$.

\subsection{Semi-infinite programs}

Now let the variable be a nonnegative measure $\Xi$ on a compact set $\mathcal{B}$, for \cref{sec:21.2} the velocities $\vect{w}$ with $|\vect{w}| \le v_{\mathrm{esc}}$ in the galactic frame. The number of constraints stays finite. Bin $i$ of an experiment predicts $\int\mathsf{k}_i\,\dd\Xi$ events, with a continuous kernel $\mathsf{k}_i(\vect{w}) \ge 0$, the count from a single stream of unit density at $\vect{w}$; the observations allow at most $N_i^{\mathrm{up}} > 0$. The quantity to bound is $\int\mathsf{f}\,\dd\Xi$ with a continuous $\mathsf{f} \ge 0$: a halo integral above some speed, or the rate in a second experiment. The program and its dual are
\begin{gather}
  \Phi_{\max} = \sup_{\Xi \ge 0}\Bigl\{\int_{\mathcal{B}}\mathsf{f}\,\dd\Xi \;:\; \int_{\mathcal{B}}\mathsf{k}_i\,\dd\Xi \le N_i^{\mathrm{up}},\ i = 1, \dots, n\Bigr\},
  \nonumber\\
  \min_{\boldsymbol{\lambda} \ge 0}\Bigl\{\sum_i\lambda_iN_i^{\mathrm{up}} \;:\; \mathsf{g}_\lambda(\vect{w}) \ge 0\ \ \forall\vect{w} \in \mathcal{B}\Bigr\},
  \label{eq:C.12-sip}
\end{gather}
with $\mathsf{g}_\lambda(\vect{w}) = \sum_i\lambda_i\mathsf{k}_i(\vect{w}) - \mathsf{f}(\vect{w})$. The primal has infinitely many variables, the values of $\Xi$ everywhere on $\mathcal{B}$; the dual has infinitely many constraints, one per point. Such programs are called \emph{semi-infinite}.

Weak duality carries over in one line. For feasible $\Xi$ and $\boldsymbol{\lambda}$,
\begin{equation}
  \int\mathsf{f}\,\dd\Xi \le \int\sum_i\lambda_i\mathsf{k}_i\,\dd\Xi = \sum_i\lambda_i\int\mathsf{k}_i\,\dd\Xi \le \sum_i\lambda_iN_i^{\mathrm{up}} ,
  \label{eq:C.12-certificate}
\end{equation}
by $\mathsf{g}_\lambda \ge 0$ and $\Xi \ge 0$, then by $\lambda_i \ge 0$ and the primal constraints. A vector $\boldsymbol{\lambda}$ with $\mathsf{g}_\lambda \ge 0$ on all of $\mathcal{B}$ is a \emph{dual certificate}: it proves $\Phi_{\max} \le \sum_i\lambda_iN_i^{\mathrm{up}}$ for every halo at once. Checking a certificate requires only the evaluation of one function on $\mathcal{B}$, which is why a halo-independent bound can be verified independently of the optimizer that produced it.

Strong duality needs hypotheses. Let $\mathcal{B}$ be compact, $\mathsf{k}_i$ and $\mathsf{f}$ continuous, and let every point where $\mathsf{f} > 0$ be seen by some bin, $\kappa_{\mathrm{v}} \equiv \min_{\vect{w} \in \operatorname{supp}\mathsf{f}}\max_i\mathsf{k}_i(\vect{w}) > 0$. Mass outside $\operatorname{supp}\mathsf{f}$ does not raise the objective and, since $\mathsf{k}_i \ge 0$, only uses up the bounds, so it can be removed. On $\operatorname{supp}\mathsf{f}$, $\sum_i\mathsf{k}_i \ge \kappa_{\mathrm{v}}$, and the constraints give
\begin{equation}
  \kappa_{\mathrm{v}}\,\Xi(\operatorname{supp}\mathsf{f}) \le \int\sum_i\mathsf{k}_i\,\dd\Xi \le \sum_iN_i^{\mathrm{up}} \equiv N_{\mathrm{tot}} ,
  \label{eq:C.12-mass}
\end{equation}
so the relevant halos have bounded total mass. The feasible set is then compact in the weak sense in which $\int\mathsf{f}\,\dd\Xi$ is continuous, and $\Xi = 0$ satisfies every constraint strictly (the Slater condition). Under these conditions the two optima of \cref{eq:C.12-sip} are equal and both are attained~\cite{shapiro2009semi}.

\subsection{Extreme points and streams}

The optimal halo is a sum of at most $n$ streams, $\Xi^* = \sum_{a=1}^{n'}\xi_a\,\delta_{\vect{w}_a}$ with $n' \le n$ and $\xi_a > 0$. On a finite grid of velocities this is the statement proved above for finite programs, with one variable per grid point. On the continuum, Rogosinski's theorem applied to the $n + 1$ functionals $\mathsf{k}_1, \dots, \mathsf{k}_n, \mathsf{f}$ replaces any optimal $\Xi$ by a measure with the same counts and objective on at most $n + 1$ points~\cite{rogosinski1958moments}, and the finite argument above, applied to those points, removes one more. Where the streams sit follows from equality in \cref{eq:C.12-certificate}, which strong duality enforces at the optimum. The first inequality is an equality only if $\Xi^*$ puts no mass where $\mathsf{g}_\lambda > 0$: the streams sit at zeros of $\mathsf{g}_\lambda$, which, since $\mathsf{g}_\lambda \ge 0$, are its minima. The second is an equality only if $\lambda_i = 0$ for every bin whose bound is not saturated. These are the \emph{complementary-slackness} conditions, and they make the optimum checkable: the streams touch the dual function from below, and only saturated bins carry weight.

With one bin the problem is solved in closed form. Then $\mathsf{g}_\lambda \ge 0$ means $\lambda_1 \ge \mathsf{f}(\vect{w})/\mathsf{k}_1(\vect{w})$ wherever $\mathsf{f} > 0$, the smallest certificate is $\lambda_1 = \sup_{\vect{w}}\mathsf{f}/\mathsf{k}_1$, and the dual value is $N_1^{\mathrm{up}}\sup_{\vect{w}}\mathsf{f}/\mathsf{k}_1$. A single stream at the maximizer $\vect{w}_*$, with $\xi_1 = N_1^{\mathrm{up}}/\mathsf{k}_1(\vect{w}_*)$, saturates the bin and achieves the same value. Hence
\begin{equation}
  \Phi_{\max} = N_1^{\mathrm{up}}\,\sup_{\vect{w} \in \mathcal{B}}\frac{\mathsf{f}(\vect{w})}{\mathsf{k}_1(\vect{w})} \qquad (n = 1) .
  \label{eq:C.12-onebin}
\end{equation}
The bound is set by the velocity at which the experiment is least sensitive relative to the quantity being bounded; every other halo does better. For a target whose response depends only on speed, $\mathsf{k}_1$ depends only on $|\vect{w}|$, and \cref{eq:C.12-onebin} becomes the familiar statement that a halo-independent bound is set by the single speed shell that the bin sees worst (\cref{sec:21.2}).

\subsection{Discretization error}

In practice the program is solved on a grid $\mathcal{B}_h \subset \mathcal{B}$ of directions and speeds, with covering radius $h$: every $\vect{w} \in \mathcal{B}$ lies within $h$ of a grid point. The grid makes the error two-sided. Restricting $\Xi$ to the grid shrinks the primal feasible set, so the grid primal optimum $\Phi_h$ obeys $\Phi_h \le \Phi_{\max}$. The grid dual solution $\boldsymbol{\lambda}_h$ satisfies $\mathsf{g}_{\lambda_h} \ge 0$ only on $\mathcal{B}_h$, and is no certificate on $\mathcal{B}$. Its violation between grid points is bounded by the Lipschitz constant $L_h$ of $\mathsf{g}_{\lambda_h}$: for any $\vect{w}$ and its nearest grid point $\vect{w}_h$, $\mathsf{g}_{\lambda_h}(\vect{w}) \ge \mathsf{g}_{\lambda_h}(\vect{w}_h) - L_h|\vect{w} - \vect{w}_h| \ge -L_hh$. Repeating \cref{eq:C.12-certificate} with $\mathsf{f} \le \sum_i\lambda_i\mathsf{k}_i + L_hh$ on $\operatorname{supp}\mathsf{f}$, and bounding the mass there with \cref{eq:C.12-mass},
\begin{equation}
  \Phi_h \;\le\; \Phi_{\max} \;\le\; \sum_i\lambda_{h,i}N_i^{\mathrm{up}} + \frac{L_h\,h\,N_{\mathrm{tot}}}{\kappa_{\mathrm{v}}} .
  \label{eq:C.12-leak}
\end{equation}
The gap between the two sides is the discretization error, and it is computable: $L_h$ follows from the derivatives of the kernels and $\kappa_{\mathrm{v}}$ from their values. It falls linearly with $h$, and faster in practice, because the optimal streams sit at minima of $\mathsf{g}_\lambda$, where the function is flat. Refining the grid only near the few streams of the grid optimum, and checking $\mathsf{g}_\lambda \ge 0$ on the refined grid, is the standard way to close it. A bound reported from a grid computation is conservative only if it uses the right side of \cref{eq:C.12-leak}.

The positivity of the kernels was used twice, to discard mass outside $\operatorname{supp}\mathsf{f}$ and in \cref{eq:C.12-mass}. If a truncated response makes some $\mathsf{k}_i$ negative at some velocities, \cref{eq:C.12-certificate} still holds for every feasible pair, but the mass bound fails, and with it the control of \cref{eq:C.12-leak}; \cref{sec:21.2} bounds that leakage separately.
